\subsection{Run-time of the Algorithm}
%\Cref{alg:WeakToStrongCover}}
\label{sec:WeakToStrongCover:strongly-poly-runtime:main}
In this section, we show that \Cref{alg:WeakToStrongCover} can be implemented to run in strongly polynomial time. 
\Cref{lem:WeakToStrong:strongly-polytime:main} is the main lemma of this section. 
In order to show that each recursive call of \Cref{alg:WeakToStrongCover} can be implemented in polynomial time, we will require the ability to compute the family \maximalTightSetFamily{p, H, w}.  
In the following lemma, we show that this family can indeed be computed in polynomial time. We defer the proof of the lemma to \Cref{appendix:sec:Function-Maximization-Oracles:computing-maximal-pHw-tight-sets-family}.

\begin{restatable}{lemma}{lemComputingMaximalTightSetFamily}\label{lem:FunctionMaximizetionOracles:computing-maximal-pHw-tight-sets-with-p-max-oracle}
 Let $p:2^V\rightarrow\Z$ be a skew-supermodular function and let $\left(H = (V, E), w:E \rightarrow\Z_+\right)$ be a hypergraph such that every hyperedge $e \in E$ is contained in some $(p, H, w)$-tight set, $b_{(H,w)}(u) \not = 0$ for every $u \in V$, and $b_{(H,w)}(X) \geq p(X)$ for every $X \subseteq V$.
Then, the family $\maximalTightSetFamily{p, H, w}$ can be computed in $O(|V|^2(|V|+|E|))$ time using $O(|V|^2)$ queries to \functionMaximizationOracle{p}. The run-time includes the time to construct the hypergraphs used as input to the queries to \functionMaximizationOracle{p}. Moreover, for each query to \functionMaximizationOracle{p}, the hypergraph $(G_0, c_0)$ used as input to the query has at most $|V|$ vertices and at most $|E| + 1$ hyperedges.  \end{restatable}

We now bound the run-time of \Cref{alg:WeakToStrongCover}. 

\begin{lemma}\label{lem:WeakToStrong:strongly-polytime:main}
Suppose that the input to \Cref{alg:WeakToStrongCover} is a non-empty hypergraph $(H = (V, E), w)$ and a symmetric skew-supermodular function $p:2^V\rightarrow \Z$ such that every hyperedge $\Tilde{e} \in E$ is contained in some $(p, H, w)$-tight set, $b_{(H,w)}(u) \not = 0$ for every $u \in V$, and $b_{(H,w)}(X) \geq p(X)$ for every $X \subseteq V$.
 Then, \Cref{alg:WeakToStrongCover}  can be implemented to run in $O(|V|^3(|V| + |E|)^2)$ time using $O(|V|^3(|V|+|E|))$ queries to \functionMaximizationOracleStrongCover{p}. The run-time includes the time to construct the hypergraphs used as input to the queries to \functionMaximizationOracleStrongCover{p}.
Moreover, for each query to \functionMaximizationOracleStrongCover{p}, the hypergraph $(G_0, c_0)$ used as input to the query has $O(|V|)$ vertices and $O(|V|+|E|)$ hyperedges.
\end{lemma}
\begin{proof}
Let $\ell$ denote the recursion depth of \Cref{alg:WeakToStrongCover}. By \Cref{lem:WeakToStrongCover:NumberOfHyperedges}, we have that $\ell=O(|E|+|V|)$. 
Let  $i \in [\ell]$ be a recursive call.
 We recall that $p_i:2^{V_i}\rightarrow\Z$ and $(H_i, w_i)$ are inputs to the $i^{th}$ recursive call, where $p_1 := p$ and $(H_1, w_1) := (H, w)$.
By Lemmas \ref{lem:weak-to-strong-cover:next-recursive-call-hypothesis-satification:symm-skew-supermod-strong-cover} and \ref{lem:weak-to-strong-cover:next-recursive-call-hypothesis-satification:hyperedge-in-tight-set-positive-coverage} and induction on $i$, we have that $p_i$ is a symmetric skew-supermodular function and $(H_i = (V_i, E_i), w_i)$ is a (non-empty if $i \leq \ell - 1$) hypergraph such that every hyperedge $\Tilde{e} \in E_i$ is contained in some $(p_i, H_i, w_i)$-tight set, $b_{(H_i, w_i)}(u) \not = 0$ for every $u \in V_i$, and $b_{(H_i, w_i)}(X) \geq p_i(X)$ for every $X \subseteq V_i$. Let $e_i, f_i, \beta_i^M, \alpha_i^M, (H_i^M, w_i^M), \beta_i^R, \alpha_i^R, (H_i^R, w_i^R)$, $p_i^R$, $(H_0^{i}, w_0^{i})$, $E_0^i, \zeros_i$ denote the quantities  as defined by \Cref{alg:WeakToStrongCover} during the $i^{th}$ recursive call. We recall that $V_i \subseteq V$. By Lemma \ref{lem:WeakToStrong:properties-of-alpha^M-and-alpha^R}(2), (5) and induction on $i$, we have that $|E_i| \leq |E| + i \leq |E| + \ell = O(|E| + |V|)$, where the final equality is by \Cref{lem:WeakToStrongCover:NumberOfHyperedges}(1). 
%Consequently, $\poly(|V_i|, |E_i|) = \poly(|V|, |E|)$ and $\poly(|V_i|) = \poly(|V|)$. 
Furthermore, the hypergraph returned in Step \ref{alg:WeakToStrongCover:(14)} contains at most $|V|$ vertices and at most $|E_i| + O(|V_i|)= O(|V| + |E|)$ hyperedges by \Cref{lem:WeakToStrongCover:NumberOfHyperedges}(2) and induction on $i$. 
The following claim shows that all steps except for Steps \ref{alg:WeakToStrongCover:(10)},  \ref{alg:WeakToStrongCover:(13)}, and  \ref{alg:WeakToStrongCover:(14)} of \Cref{alg:WeakToStrongCover} in the $i^{th}$ recursive call can be implemented to run in $O(|V_i|^2(|V_i| + |E_i|))$ time
using $O(|V_i|^2)$ queries \functionMaximizationOracle{p_i}. This includes the run-time to construct the hypergraphs used in the queries to \functionMaximizationOracle{p_i}. We will address the run-time for Steps \ref{alg:WeakToStrongCover:(10)},  \ref{alg:WeakToStrongCover:(13)}, and  \ref{alg:WeakToStrongCover:(14)} later. 

\begin{claim}\label{claim:weak-to-strong-runtime-of-one-recursive-call}
For every $i\in [\ell]$, all steps except for Steps \ref{alg:WeakToStrongCover:(10)},  \ref{alg:WeakToStrongCover:(13)}, and  \ref{alg:WeakToStrongCover:(14)} of \Cref{alg:WeakToStrongCover} in the $i^{th}$ recursive call can be implemented to run in 
$O(|V_i|^2(|V_i| + |E_i|))$ time
using 
$O(|V_i|^2)$ queries to \functionMaximizationOracle{p_i}. The run-time includes the time to construct the hypergraphs used as input to the queries to \functionMaximizationOracle{p_i}. Moreover, for each query to \functionMaximizationOracle{p_i}, the hypergraph $(G_0, c_0)$ used as input to the query has  $|V_i|$ vertices and at most $|E_i| + 1$ hyperedges.
\end{claim}
\begin{proof}
We first show how to compute hyperedges $e_i$ and $f_i$ (Step \ref{alg:WeakToStrongCover:(4)}). We recall that by \Cref{cor:weak-to-strong:uncrossing:exist-two-hyperedges-contained-in-disjoint-maximal-tight-sets}, there exist two hyperedges $e_i, f_i \in E_i$ such that $\maximalTightSetContainingHyperedge{p_i, H_i, w_i}(e_i) \not = \maximalTightSetContainingHyperedge{p_i, H_i, w_i}(f_i)$ and consequently Step \ref{alg:WeakToStrongCover:(4)} is well defined. By \Cref{lem:weak-to-strong:uncrossing:maximal-tight-set-family-disjoint}, the family $\maximalTightSetFamily{p_i, H_i, w_i}$ is a disjoint family. 
%Thus, $|\maximalTightSetFamily{p_i, H_i, w_i}| = O(|V_i|)$. 
Furthermore,  by \Cref{lem:FunctionMaximizetionOracles:computing-maximal-pHw-tight-sets-with-p-max-oracle}, we can compute the family $\maximalTightSetFamily{p_i, H_i, w_i}$ in $O(|V_i|^2(|V_i|+|E_i|))$ time using $O(|V_i|^2)$ queries to \functionMaximizationOracle{p_i}. 
Moreover, for each query to \functionMaximizationOracle{p_i}, the hypergraph $(G_0, c_0)$ used as input to the query has at most $|V_i|$ vertices and at most $|E_i|+1$ hyperedges. 
We note that $\maximalTightSetContainingHyperedge{p_i, H_i, w_i}(\Tilde{e}) \in \maximalTightSetFamily{p_i, H_i, w_i}$ for every $\Tilde{e} \in E$.
Consequently, we can find the required hyperedges for Step \ref{alg:WeakToStrongCover:(4)} by finding a pair of distinct sets $X, Y \in \maximalTightSetFamily{p_i, H_i, w_i}$ that each contains a hyperedge, and define these hyperedges to be $e_i$ and $f_i$ respectively. Thus, Step  \ref{alg:WeakToStrongCover:(4)} can be implemented to run in $O(|V_i|^2(|V_i|+|E_i|))$ time using $O(|V_i|^2)$ queries to \functionMaximizationOracle{p_i}. Moreover, for each query to \functionMaximizationOracle{p_i}, the hypergraph $(G_0, c_0)$ used as input to the query has at most $|V_i|$ vertices and at most $|E_i|+1$ hyperedges. 

Next, we show how to compute the values $\alpha_i^M$ (Step \ref{alg:WeakToStrongCover:(5)}) and $\alpha^R_i$ (Step \ref{alg:WeakToStrongCover:(7)}).  We first show how to compute the value $\beta_i^M$.
We have the following:
\begin{align*}
    \beta_i^M& = \min\left\{b_{(H_i,w_i)}(X) - p_i(X) : X \subseteq V_i \text{ and } e_i \cap X, f_i \cap X \not = \emptyset\right\} &\\
    & = \min_{u \in e_i, v \in f_i}\left\{\min\left\{b_{(H_i,w_i)}(X) - p_i(X) : \{u,v\} \subseteq X \subseteq V_i\right\}\right\}.&
    %&=: \min_{u \in e_i, v \in f_i}\left\{\beta_i^{uv}\right\},&
\end{align*}
For $u \in e_i$ and $v \in f_i$, let $\beta_i^{uv} := \min\{b_{(H_i,w_i)}(X) - p_i(X) : \{u,v\} \subseteq X \subseteq V_i\}.$ Thus, $\beta_i^M = \min_{u \in e_i, v\in f_i}\beta_i^{uv}$.
Then, $\beta_i^M$ (and consequently the value $\alpha_i^M$) can be computed in $O(|V_i|^2(|V_i| + |E_i|))$ time using 
$|V_i|^2$ queries to \functionMaximizationOracle{p_i}, where each query is on a input hypergraph $(G_0, c_0)$ that has at most $|V_i|$ vertices and at most $|E_i|+1$ hyperedges. This is because for every $u \in e_i$ and $v \in f_i$ (there are $O(|V_i|^2)$ choices for $u$ and $v$), the value $\beta_i^{uv}$ can be computed using a single oracle call to \functionMaximizationOracle{p_i} on a hypergraph $(G_0, c_0)$ that has at most $|V_i|$ vertices and at most $|E_i|+1$ hyperedges.
Next, the value $\beta^R_i$, (and consequently the value $\alpha^R_i$)  can be computed 
in $O(|V_i| + |E_i|)$ time using a single query to \functionMaximizationOracle{p_i} on the input hypergraph $(G_0=H_i^M, c_0=w_i^M)$ that has  $|V_i|$ vertices and at most $|E_i|+1$ hyperedges. Here, the hypergraph $(H_i^M, w_i^M)$ has  $|V_i|$ vertices and $|E_i|+1$ hyperedges because of  \Cref{lem:WeakToStrong:properties-of-alpha^M-and-alpha^R}(2).
Next, we show that the hypergraphs $(H_i^M, w_i^M)$ (Step \ref{alg:WeakToStrongCover:(6)}) and $(H_i^R, w_i^R)$ (Step \ref{alg:WeakToStrongCover:(8)}) can be computed in $O(|V_i| + |E_i|)$ time. We recall that the \merge operation decreases the weights of hyperedges $e_i$ and $f_i$ by $\alpha_i^M$,  increases the weight of the hyperedge $e_i\cup f_i$ by $\alpha_i^M$ and discards the resulting zero-weighted edges (if any). This can be implemented to run in $O(|E_i|)$ time. Furthermore, the \reduce operation decreases the weight of the hyperedge $e_i\cup f_i$ by $\alpha^R_i$, and so can also be implemented to run in $O(|E^R_i|) = O(|E_i|)$ time by \Cref{lem:WeakToStrong:properties-of-alpha^M-and-alpha^R}(2) and (5). Finally, the computations of the hypergraph $(H^{i}_0, w^{i}_0)$ (Step \ref{alg:WeakToStrongCover:(9)}), the set $\zeros_i$ (\Cref{alg:WeakToStrongCover:(11)}) and the set $V'_i$  (\Cref{alg:WeakToStrongCover:(12)}) can each be done in $O(|V_i|)$ time. Furthermore, the hypergraph returned in Step \ref{alg:WeakToStrongCover:(14)} contains at most $|V|$ vertices and at most $|E_i| + O(|V_i|)$ hyperedges. Consequently, Steps \ref{alg:WeakToStrongCover:(15)} and \ref{alg:WeakToStrongCover:(16)}  can be implemented to run in $O(|V_i| + |E_i|)$ time. 
\end{proof}

\Cref{claim:weak-to-strong-runtime-of-one-recursive-call} implies that each query to \functionMaximizationOracle{p_i} takes as input a hypergraph 
with at most $|V_i|$ vertices and $O(|E_i|)$ hyperedges.  
Next, we focus on the time to implement a query to \functionMaximizationOracle{p_i} on such hypergraphs. 

\begin{claim}\label{claim:WeakToStrong:p_i-oracle-using-p-oracle}
For every $i\in [\ell]$ and for every given hypergraph $(G = (V_i, E_G), c:E_G\rightarrow\Z_+)$ such that $|E_G| = O(|E_i|)$ AND disjoint subsets $S, T\subseteq V_i$, $\functionMaximizationOracle{p_i}((G, c), S, T, y)$ can be computed in $O(|V|(|V| + |E|))$ time using at most $|V_i|+1$ queries to \functionMaximizationOracleStrongCover{p}. The run-time includes the time to construct the hypergraphs used as input to the queries to \functionMaximizationOracleStrongCover{p}. Moreover, for each query to \functionMaximizationOracleStrongCover{p}, the hypergraph $(G_0, c_0)$ used as input to the query has $O(|V|)$ vertices and $O(|V|+|E|)$ hyperedges.
\end{claim}
\begin{proof}
Let the input to \functionMaximizationOracle{p_i} be a hypergraph $(G = (V_i, E_G), c:E_G\rightarrow\Z_+)$ such that $|E_G| = O(|E_i|)$ and disjoint subsets $S, T\subseteq V_i$. By \Cref{lem:Preliminaries:wc-oracle-from-sc-oracle}, we can implement $\functionMaximizationOracle{p_i}((G, c), S, T, y_0)$ in $O(|V_i|(|V_i| + |E_i|))$ time using at most $|V_i|+1$ queries to $\functionMaximizationOracleStrongCover{p_i}$, where each query to $\functionMaximizationOracleStrongCover{p_i}$ is on a input hypergraph with $|V_i|+1$ vertices and $O(|E_i|)$ hyperedges. Next, we show that the query $\functionMaximizationOracleStrongCover{p_i}((G, c), S, T)$ where $(G = (V_i, E_G), c:E_G\rightarrow\Z_+)$ is a hypergraph with $|E_G| = O(|E_i|)$ and $S, T\subseteq V_i$ are disjoint subsets, can be implemented in $O(|V| + |E|)$ time using one query to  $\functionMaximizationOracleStrongCover{p}$ on an input hypergraph $(G_0, c_0)$ that has $O(|V|)$ vertices and $O(|V|+|E|)$ hyperedges. 

We first recall some notation and define some quantities for convenience. For $i\in [\ell]$, we recall that $(H_0^i, w_0^i)$ is the hypergraph $(H_0, w_0)$, $p^R_i, p_i'$ are the functions $p^R, p'$ respectively, and $\zeros_i$ is the set $\zeros$ defined in the $i^{th}$ recursive call of the algorithm. We note that $p_i' =p_{i+1}$ for every $i\in [\ell-1]$. Furthermore, the hypergraph $(H_0^i, w_0^i)$ contains exactly one hyperedge for every $i\in [\ell-1]$. 
For $i\in [\ell-1]$, let $(H_1^i, w_1^i)$ be the hypergraph obtained from $(H_0^i, w_0^i)$ by adding the vertices $\cup_{j=1}^{i-1}\zeros_{j}$.

Let $i\in [\ell]$. Let the input to \functionMaximizationOracleStrongCover{p_i} be a hypergraph $(G = (V_i, E_G), c:E_G\rightarrow\Z_+)$ such that $|E_G| = O(|E_i|)$ and disjoint subsets $S_0, T_0\subseteq V_i$. 
We note that $p_i=p_{i-1}'$ and hence, it suffices to answer $\functionMaximizationOracleStrongCover{p_{i-1}'}((G, c), S_0, T_0)$. 
Construct the hypergraph $(G_1, c_1)$ to be the hypergraph obtained from $(G,c)$ by adding vertices $\cup_{j=1}^{i-1}\zeros_{j}$. 
Construct the hypergraph $(\Tilde{H}, \Tilde{w})=(\sum_{j=1}^{i-1}H_1^j, \sum_{j=1}^i w_1^j)$. The hypergraph $(\Tilde{H}, \Tilde{w})$ contains at most $i - 1 \leq \ell = O(|E| + |V|)$ hyperedges, and can be constructed in time $O(|V|+|E|)$. 
Next, construct the hypergraph $(G, c)=(G_1+\Tilde{H}, c_1+\Tilde{w})$. The hypergraph $(G, c)$ contains $O(|E|+|V|)$ hyperedges and can be constructed in time $O(|V|+|E|)$. 
%Let $y_1\in \R^{V}$ be the vector obtained by setting $y_1(u):=y_0(u)$ for every $u\in V_{i}$ and $y_1(u):=0$ for every $u\in V\setminus V_{i}$. We note that the vector $y_1$ can be constructed in $O(|V|)$ time.  
Obtain $(Z, p(Z))$ by querying $\functionMaximizationOracleStrongCover{p}((G, c), S_0, T_0)$ and return $(Z-\cup_{j=1}^{i-1}\zeros_{j}, p(Z))$. 
\end{proof}

By \Cref{claim:weak-to-strong-runtime-of-one-recursive-call}, each query to \functionMaximizationOracle{p_i} takes as input a hypergraph with $|V_i|\le |V|$ vertices and $O(|E_i|) = O(|E|+|V|)$ hyperedges. 
By \Cref{claim:WeakToStrong:p_i-oracle-using-p-oracle}, a query to \functionMaximizationOracle{p_i} on a input hypergraph with at most $|V_i|$ vertices and at most $O(|E_i|)$ hyperedges 
can be implemented in $O(|V|(|V| + |E|))$  time using $O(|V|)$ queries to \functionMaximizationOracleStrongCover{p}, where each query to \functionMaximizationOracleStrongCover{p} is on a input hypergraph $(G_0, c_0)$ that has $O(|V|)$ vertices and $O(|E|+|V|)$ hyperedges. 
Thus, all steps except for Steps \ref{alg:WeakToStrongCover:(10)}, \ref{alg:WeakToStrongCover:(13)}, and \ref{alg:WeakToStrongCover:(14)} of \Cref{alg:WeakToStrongCover} in the $i^{th}$ recursive call can be implemented to run in $O(|V|^3(|V| + |E|))$ time using $O(|V|^3)$ queries to \functionMaximizationOracleStrongCover{p}, where each query to \functionMaximizationOracleStrongCover{p} is on a input hypergraph $(G_0, c_0)$ that has $O(|V|)$ vertices and $O(|E|+|V|)$ hyperedges. 

We note that Steps \ref{alg:WeakToStrongCover:(10)} and \ref{alg:WeakToStrongCover:(13)} need not be implemented explicitly for the purposes of the algorithm. Instead, \functionMaximizationOracle{p_{i+1}} on input hypergraphs with $|V_{i+1}|\le |V|$ vertices and at most $|E_{i+1}|+1$ hyperedges can be used to execute all steps except for Steps \ref{alg:WeakToStrongCover:(10)}, \ref{alg:WeakToStrongCover:(13)}, and \ref{alg:WeakToStrongCover:(14)} in the $(i+1)^{th}$ recursive call. 
We have already seen in \Cref{claim:WeakToStrong:p_i-oracle-using-p-oracle} that a query to \functionMaximizationOracle{p_{i+1}} on input hypergraphs with at most $|V|$ vertices and $O(|E_{i+1}|) = O(|V| + |E|)$ hyperedges can be implemented in  $O(|V|(|V|+|E|))$ time using $O(|V|)$ queries to \functionMaximizationOracleStrongCover{p}. 
Step \ref{alg:WeakToStrongCover:(14)} is the recursive call. We note that the hypergraph returned in Step \ref{alg:WeakToStrongCover:(14)} contains at most $|V|$ vertices and at most $|E_i| + O(|V_i|)= O(|V| + |E|)$ hyperedges. 

Thus, we have shown that each recursive call can be implemented to run in time $O(|V|^3(|V| + |E|))$ using $O(|V|^3)$ queries to \functionMaximizationOracleStrongCover{p}, where each query to \functionMaximizationOracleStrongCover{p} is on an input hypergraph $(G_0, c_0)$ that has $O(|V|)$ vertices and $O(|E|+|V|)$ hyperedges. The number of recursive calls is $O(|V|+|E|)$. 
Hence, the algorithm can be implemented to run in $O(|V|^3(|V| + |E|)^2)$ time using $O(|V|^3(|V|+|E|))$ queries to \functionMaximizationOracleStrongCover{p}, where each query to \functionMaximizationOracleStrongCover{p} is on a input hypergraph $(G_0, c_0)$ that has $O(|V|)$ vertices and $O(|E|+|V|)$ hyperedges. 
\end{proof}